\onecolumn
\begin{tcolorbox}[
    title=Training Data Format of \method: Correction Trajectories,
    colback=white,
    colframe=black!75!white,
    breakable,
    boxrule=0.8pt,
    left=2mm,
    right=2mm,
    top=1mm,
    bottom=1mm
]
\refstepcounter{tcbtable}
\label{tab:traing_data}
\begin{spacing}{1.1}

\textbf{/**Input**/}  

1. People numbered $1,2,\cdots,2015$ are arranged in a line, and a position-swapping game is played among them, with the rule that each swap can only occur between adjacent individuals. Now, the person numbered 100 and the person numbered 1000 are to swap positions, with the minimum number of swaps required being $\qquad$ times.  

\textbf{/**Output**/}  

\textbf{\color{purple}{<answer>}}  

Let's reason through this problem:

We need to find the minimum number of swaps needed for person 100 and person 1000 to swap their positions.

The key here is understanding that every swap moves two people forward or backward by one place. Specifically, we want to move person 100 and person 1000 relative to each other by 900 positions (from 100 to 1000). Since each swap affects two positions at once, 

\vspace{\medskipamount}
\textit{\color{gray}{( ... more reasoning tokens omitted for brevity )}} 

Considering efficiency and the simplicity of the problem, we calculate:

\[
\text{Minimum swaps required} = 900 + (\text{Efficiency savings})
\]

Given our targeted statement on "minimum", we assume efficiency yields ideal movement, thus aligning our answer closely with Calculation 1:

\[
\boxed{900}
\]
\textbf{\color{purple}{</answer>}}  

\textbf{\color{purple}{</critic>}}  

The key mistake is that the answer incorrectly equates the number of swaps with the difference in the initial positions of the two individuals.

This is wrong because it only accounts for the movement of one person to the other's position, ignoring the movement required for the second person. To swap person A (at position 100) and person B (at position 1000), 
\vspace{\medskipamount}
\textit{\color{gray}{( ... more reasoning tokens omitted for brevity )}} 
% both must traverse the space between them. A must move past the 899 people in between and then B (900 swaps). This leaves B at position 999. B must then move past those same 899 people to get to position 100 (899 more swaps).

\textbf{F.}

\textbf{\color{purple}{</critic>}}  

\textbf{\color{purple}{<revised>}}  

To improve the answer, it’s important to recognize the full process for swapping two individuals who are not adjacent.

\vspace{\medskipamount}
\textit{\color{gray}{( ... more reasoning tokens omitted for brevity )}} 

This is a distance of 899 positions, requiring 899 swaps. Each swap moves person 1000 one spot towards position 100.

Thus, the correct minimum number of adjacent swaps required is the sum:

\[
900 \text{ (to move 100 to 1000)} + 899 \text{ (to move 1000 to 100)} = 1799
\]

\[
\boxed{1799}
\]

\textbf{\color{purple}{<revised>}}  
\end{spacing}

\end{tcolorbox}
